\DeclareEditionOverlay{CM-C06-U001}{%
\begin{EditionUnitCard}{La idea local}
Compartir derivadas en el centro controla el contacto cerca de ese punto. El
orden no promete igualdad lejos del centro ni fija siempre el grado efectivo.
\end{EditionUnitCard}}

\DeclareEditionOverlay{CM-C06-U002}{%
\begin{EditionUnitCard}{De las derivadas a los coeficientes}
Sigue la construcción en dos direcciones: qué exige cada derivada en el centro
y por qué esas exigencias determinan un único coeficiente.
\end{EditionUnitCard}}

\DeclareEditionOverlay{CM-C06-U003}{%
\Needspace{0.72\textheight}
\begin{EditionUnitCard}{Comparar sin perder el centro}
\RegisteredBookFigure{FIG-C06-001}{fig:c06:taylor-approximation}{Comparación de
$\sin x$ con sus polinomios $P_1$, $P_3$ y $P_5$.}{Gráfico cartesiano en el
intervalo de menos pi a pi: la curva seno se compara con tres polinomios de
Taylor; cerca de cero las curvas de mayor grado permanecen próximas durante un
intervalo más amplio.}{\TaylorApproximationVisual}
\begin{EditionCheckpointBox}{Piénsalo}
Elige un punto cercano a cero y otro cercano a $\pi$. ¿Qué cambia al comparar
los tres grados? Formula la respuesta como observación local, no global.
\end{EditionCheckpointBox}
\begin{EditionWarningBox}{Polinomio finito, no serie infinita}
Cada curva mostrada corresponde a una suma finita. El gráfico no afirma que
toda función suave coincida con una serie de Taylor.
\end{EditionWarningBox}
\end{EditionUnitCard}}

\DeclareEditionOverlay{CM-C06-U004}{%
\begin{EditionUnitCard}{Leer el teorema antes de usarlo}
Separa hipótesis sobre el intervalo, continuidad de derivadas y existencia de
la derivada siguiente. Después identifica qué garantiza el punto intermedio.
\begin{EditionWarningBox}{El punto intermedio no es un dato}
El teorema garantiza que existe; normalmente no se calcula ni se mantiene fijo
al cambiar $x$.
\end{EditionWarningBox}
\end{EditionUnitCard}}

\DeclareEditionOverlay{CM-C06-U005}{%
\Needspace{0.72\textheight}
\begin{EditionUnitCard}{Del resto a una garantía}
\RegisteredBookFigure{FIG-C06-002}{fig:c06:taylor-error}{Error absoluto de
$P_1$, $P_3$ y $P_5$ al aproximar $\sin x$.}{Tres curvas de error absoluto en
el intervalo de menos pi a pi se anulan en cero y crecen al alejarse; los
estilos de línea distinguen cada grado sin depender del color.}{\TaylorErrorVisual}
\begin{EditionCheckpointBox}{Antes de seguir}
Fija un mismo valor de $x$ y ordena los tres errores observados. ¿Por qué ese
orden visual no sustituye una cota garantizada para todo un intervalo?
\end{EditionCheckpointBox}
\begin{EditionWarningBox}{La cota cubre un segmento}
No basta conocer la derivada en el centro: el número $M$ debe controlarla en
todo el tramo relevante.
\end{EditionWarningBox}
\end{EditionUnitCard}}

\DeclareEditionOverlay{CM-C06-U006}{%
\begin{EditionUnitCard}{La familia importa}
En seno y coseno las derivadas se acotan de forma uniforme. En el logaritmo,
el intervalo no puede cruzar la frontera del dominio. Conserva esa diferencia
cuando leas los tres casos.
\end{EditionUnitCard}}

\DeclareEditionOverlay{CM-C06-U007}{%
\begin{EditionUnitCard}{Aproximación y certificado}
Anota por separado el valor producido por el polinomio y el máximo error que
permite afirmar el teorema. Una cifra decimal no queda justificada solo porque
aparezca en la calculadora.
\end{EditionUnitCard}}

\DeclareEditionOverlay{CM-C06-U008}{%
\begin{EditionUnitCard}{El signo también informa}
Una forma de resto puede servir para algo más que medir tamaño: su signo puede
convertir la aproximación en una desigualdad. Revisa siempre el intervalo que
hace válido ese argumento.
\end{EditionUnitCard}}

\DeclareEditionOverlay{CM-C06-U009}{%
\begin{EditionUnitCard}{Elegir el orden}
Cuando hay cancelaciones, conserva términos hasta encontrar el primero que no
se anula. Después pregunta si Taylor es realmente el método más directo.
\end{EditionUnitCard}}

\DeclareEditionOverlay{CM-C06-U010}{%
\begin{EditionUnitCard}{Componer requiere datos adicionales}
En una sustitución con $o$ pequeña, identifica la variable que tiende, el
centro, el orden y dónde está definida la composición. Para volver a potencias
de $x-a$ necesitas conocer además el orden de la función interior.
\end{EditionUnitCard}}

\DeclareEditionOverlay{CM-C06-U011}{%
\begin{EditionUnitCard}{Lectura de ampliación}
La forma de Cauchy amplía el repertorio, pero no desplaza el resto de Lagrange.
El resto integral queda fuera porque depende de integración, no porque sea un
resultado falso.
\end{EditionUnitCard}}

\DeclareEditionOverlay{CM-C06-U012}{%
\begin{EditionCheckpointBox}{Recupera la distinción}
Para justificar una precisión decimal en un intervalo, ¿citarías una curva de
error observada o una cota del resto? Explica qué garantía aporta tu elección.
\end{EditionCheckpointBox}}

